\documentclass[11pt]{article}
\usepackage[margin=1in]{geometry}
\usepackage{amsmath, amssymb}
\usepackage{xcolor}
\usepackage{mdframed}
\usepackage{parskip}
\usepackage{enumitem}
\usepackage{array}
\usepackage{booktabs}

\newmdenv[
  linecolor=black!25,
  backgroundcolor=black!3,
  roundcorner=4pt,
  skipabove=8pt,
  skipbelow=8pt,
  innertopmargin=6pt,
  innerbottommargin=6pt
]{keybox}

\title{The 128-Dimensional SIFT Descriptor}
\author{}
\date{}

\begin{document}
\maketitle

\section{Setup: which image, and the one input you already have}

A keypoint carries $(s_{\text{idx}}, y, x, \sigma)$ --- one row of the $(N,4)$ tensor. The
descriptor is built from the \textbf{spatial gradients of a single Gaussian-blurred image}:
the slice chosen by matching the keypoint's $\sigma$ against the octave's $\sigma$ ladder.
\emph{Not} the DoG, and \emph{not} the 3-D scale volume.

Once that slice is picked, \textbf{everything below is 2-D image space}. Scale has only two
remaining jobs: it selected the slice, and (through $\sigma$) it sets the size of the
sampling window.

\section{Step 1 --- Gradients as magnitude and angle}

On the chosen blurred image $L$, at each pixel, from the two axis-aligned central
differences:
\[
G_x = L_{x+1} - L_{x-1}, \qquad G_y = L_{y+1} - L_{y-1},
\]
\[
m = \sqrt{G_x^2 + G_y^2}, \qquad \theta = \operatorname{atan2}(G_y,\ G_x).
\]
The magnitude $m$ is edge strength; the angle $\theta$ is edge direction. The angle is
\emph{measured} from the two differences --- never requested as an input.

\section{Step 2 --- Dominant orientation (rotation invariance)}

Before building anything, find the keypoint's own ``up'':
\begin{itemize}[nosep]
  \item Over a region around the keypoint (radius $\propto \sigma$), build a
        \textbf{36-bin} histogram ($10^\circ$ each) of $\theta$.
  \item Each pixel votes by its magnitude $m$, Gaussian-weighted so nearer pixels count more.
  \item The \textbf{peak bin} is the keypoint's reference orientation.
\end{itemize}
Every angle used afterward is measured \emph{relative} to this reference:
\[
\theta_{\text{relative}} = \theta_{\text{measured}} - \theta_{\text{dominant}}.
\]
\begin{keybox}
Rotation invariance comes entirely from this subtraction: rotate the frame to the feature's
own dominant direction, so the same feature photographed rotated yields the same descriptor.
The image is still only ever differenced along $x$ and $y$ --- the rotation is arithmetic on
the angle, not a tilted derivative.
\end{keybox}
If a second histogram peak reaches $\ge 80\%$ of the maximum, a \textbf{second keypoint} is
emitted at the same $(x, y, \sigma)$ with that orientation --- so one detection can produce
two descriptors (descriptor count $\ge$ keypoint count).

\section{Step 3 --- A $4\times 4$ grid of cells}

Center a window on the keypoint, \textbf{rotated} to the dominant orientation and sized
$\propto \sigma$, and split it into a $4\times 4$ arrangement of subregions:
\[
4 \times 4 = 16 \text{ cells}.
\]
Because the grid is rotated, its sample points fall at \emph{fractional} image locations, so
$L$ is read there by \textbf{bilinear interpolation}.

\section{Step 4 --- An 8-bin orientation histogram per cell}

Within each of the 16 cells, build an \textbf{8-bin} histogram ($45^\circ$ each) of the
gradient orientations there --- each pixel voting by $m$, Gaussian-weighted from the center,
all angles \emph{relative to the dominant orientation}.

\section{Step 5 --- That is the 128}

\begin{keybox}
\[
\underbrace{4 \times 4}_{\text{spatial cells}} \;\times\; \underbrace{8}_{\text{orientation bins}}
\;=\; 128.
\]
\end{keybox}
Flatten the 16 per-cell histograms into one vector of length 128.

\section{Step 6 --- Normalize (illumination invariance)}

\begin{enumerate}[nosep]
  \item Normalize the 128-vector to unit length --- cancels contrast / gain (multiplicative
        light changes).
  \item Clamp every element to $\le 0.2$ --- stops one huge gradient from dominating
        (robustness to non-linear illumination, e.g.\ glare).
  \item Renormalize to unit length.
\end{enumerate}
The result is the final descriptor.

\section{The interpolation that separates working SIFT from fragile SIFT}

Three places, all so the descriptor changes \emph{smoothly} as the keypoint shifts slightly
--- essential for matching across viewpoints:
\begin{itemize}[nosep]
  \item \textbf{Rotated window} $\rightarrow$ bilinear image interpolation at fractional
        sample points (Step 3).
  \item \textbf{Trilinear vote spreading:} each gradient sample is split across neighbouring
        cells (2 spatial axes) and neighbouring orientation bins (1 angular axis), weighted
        by distance to each bin center, so a pixel near a boundary does not jump
        discontinuously between cells. \emph{This is the single biggest contributor to
        descriptor stability.}
  \item \textbf{$\sigma$ scales the window} (Step 3) --- the only place scale re-enters, just
        a multiplier on the sampling radius.
\end{itemize}

\section{Why this design gives the invariances}

\begin{center}
\renewcommand{\arraystretch}{1.3}
\begin{tabular}{>{\raggedright}p{0.32\textwidth} p{0.58\textwidth}}
\toprule
\textbf{Property} & \textbf{Where it comes from} \\
\midrule
Scale invariance        & descriptor built in the keypoint's octave, window sized by $\sigma$ \\
Rotation invariance     & all angles relative to the dominant orientation (Step 2) \\
Illumination invariance & unit-normalize $\rightarrow$ clamp $0.2$ $\rightarrow$ renormalize (Step 6) \\
Small-shift robustness  & Gaussian weighting + trilinear vote spreading (Step 4) \\
\bottomrule
\end{tabular}
\end{center}

That combination --- scale-, rotation-, and light-invariant, and stable under small position
error --- is what lets the \emph{same} descriptor match across two photos taken from
different distances and angles: the whole reason the descriptors are computed for
reconstruction.

\end{document}
